\section{AHSME 1960}
        \begin{enumerate}
        	\item (\textit{AHSME 1960}) If $2$ is a solution (root) of $x^3+hx+10=0$, then $h$ equals:


 $\textbf{(A)}10\qquad \textbf{(B)}9 \qquad \textbf{(C)}2\qquad \textbf{(D)}-2\qquad \textbf{(E)}-9$


 
 \textbf{Answer: }E



\textbf{Solution}

Substitute $2$ for $x$. We are given that this equation is true. Thus, 

 $2^3+2h+10 =0$

 $18+2h=0$

 $2h=-18$

 $h=-9$

 Thus, the answer is $\boxed{\textbf{(E) }-9}$.

 


	\item (\textit{AHSME 1960}) It takes $5$ seconds for a clock to strike $6$ o'clock beginning at $6{:}00$ o'clock precisely. 
If the strikings are uniformly spaced, how long, in seconds, does it take to strike $12$ o'clock?


 $\textbf{(A)}9\frac{1}{5}\qquad \textbf{(B )}10\qquad \textbf{(C )}11\qquad \textbf{(D )}14\frac{2}{5}\qquad \textbf{(E )}\text{none of these}$


 
 \textbf{Answer: }C



\textbf{Solution}

Between six strikes, there are five intervals of space. Since it takes five seconds total, each interval is one second long. Therefore, $12$ strikes will take $11$ seconds. $\boxed{\textbf{(C)}}$.

 


	\item (\textit{AHSME 1960}) Applied to a bill for $\textdollar{10,000}$ the difference between a discount of $40$\% and two successive discounts of $36$\% and $4$\%, 
expressed in dollars, is:


 $\textbf{(A) }0\qquad \textbf{(B) }144\qquad \textbf{(C) }256\qquad \textbf{(D) }400\qquad \textbf{(E) }416$


 
 \textbf{Answer: }B



\textbf{Solution}

Taking the discount of $40$\% means you're only paying $60$\% of the bill.  That results in $10,000\cdot0.6=\textdollar{6,000}$.

 Likewise, taking two discounts of $36$\% and $4$\% means taking $64$\% of the original amount and then $96$\% of the result.  That results in $10,000\cdot0.64\cdot0.96=\textdollar{6,144}$.

 Taking the difference results in $6,144-6,000=\textdollar{144}$, or answer choice $\boxed{\textbf{(B)}}$.

 


	\item (\textit{AHSME 1960}) Each of two angles of a triangle is $60^{\circ}$ and the included side is $4$ inches. The area of the triangle, in square inches, is:


 $\textbf{(A) }8\sqrt{3}\qquad \textbf{(B) }8\qquad \textbf{(C) }4\sqrt{3}\qquad \textbf{(D) }4\qquad \textbf{(E) }2\sqrt{3}$


 
 \textbf{Answer: }C



\textbf{Solution}

If two of the angles are $60^{\circ}$, then the other angle is $60^{\circ}$ because angles in triangle add up to $180^{\circ}$.  That makes the triangle an \href{https://artofproblemsolving.com/wiki/index.php?title=Equilateral\_triangle}{equilateral triangle}, so all sides are $4$ inches long.

 
                    \begin{figure}[H]
                        centering
                        \includegraphics[width=0.5\textwidth]{Solution/AHSME_I_1960_4_Solution_Figure_1.png}
                    \end{figure}
                    

 Using the area formula $A = \frac{s^2\sqrt{3}}{4}$, the area of the triangle is $\frac{4^2\sqrt{3}}{4} = 4\sqrt{3}$ square inches, which is answer choice $\boxed{\textbf{(C)}}$.

 


	\item (\textit{AHSME 1960}) The number of distinct points common to the graphs of $x^2+y^2=9$ and $y^2=9$ is:


 $\textbf{(A) }\text{infinitely many}\qquad \textbf{(B) }\text{four}\qquad \textbf{(C) }\text{two}\qquad \textbf{(D) }\text{one}\qquad \textbf{(E) }\text{none}$


 
 \textbf{Answer: }C



\textbf{Solution}

Solve the second equation by taking the square root of both sides.
 \[y^2=9\]
 \[y=\pm3\]

 Solve the first equation by substituting $y^2$ into the first equation then solving for $x$.
 \[x^2+9=9\]
 \[x^2=0\]
 \[x=0\]

 The two solutions are $(0,3)$ and $(0,-3)$, so the answer is $\boxed{\textbf{(C)}}$.

 


	\item (\textit{AHSME 1960}) The circumference of a circle is $100$ inches. The side of a square inscribed in this circle, expressed in inches, is:


 $\textbf{(A) }\frac{25\sqrt{2}}{\pi}\qquad \textbf{(B) }\frac{50\sqrt{2}}{\pi}\qquad \textbf{(C) }\frac{100}{\pi}\qquad \textbf{(D) }\frac{100\sqrt{2}}{\pi}\qquad \textbf{(E) }50\sqrt{2}$


 
 \textbf{Answer: }B



\textbf{Solution}

First, find the diameter of the circle.  Plug in the circumference for the formula $C = \pi d$ to solve for $d$.
 \[100 = \pi d\]
 \[d = \frac{100}{\pi}\]Since all of an inscribed square's vertices touch the circle, and one of the angles of the circle is $90^{\circ}$, the square's diagonal is the diameter of the circle.

 
                    \begin{figure}[H]
                        centering
                        \includegraphics[width=0.5\textwidth]{Solution/AHSME_I_1960_6_Solution_Figure_1.png}
                    \end{figure}
                    

 From the Pythagorean theorem (or 45-45-90 triangles), the side length is $\frac{50\sqrt{2}}{\pi}$ inches, which is answer choice $\boxed{\textbf{(B)}}$ .

 


	\item (\textit{AHSME 1960}) Circle $I$ passes through the center of, and is tangent to, circle $II$. The area of circle $I$ is $4$ square inches. 
Then the area of circle $II$, in square inches, is:


 $\textbf{(A) }8\qquad \textbf{(B) }8\sqrt{2}\qquad \textbf{(C) }8\sqrt{\pi}\qquad \textbf{(D) }16\qquad \textbf{(E) }16\sqrt{2}$


 
 \textbf{Answer: }D



\textbf{Solutions}


                    \begin{figure}[H]
                        centering
                        \includegraphics[width=0.5\textwidth]{Solution/AHSME_I_1960_7_Solution_Figure_1.png}
                    \end{figure}
                    

 



\textbf{Solution 1}

Since Circle $I$ is tangent to circle $II$ and touches the center of circle $II$, the diameter of circle $I$ is the radius of circle $II$.

 That means circle $II$ is twice as big as circle $I$, so the area of circle $II$ is four times as big as circle $I$.

 The area of circle $II$ is $16$ square inches, so the answer is $\boxed{\textbf{(D)}}$.

 



\textbf{Solution 2}

Since Circle $I$ is tangent to circle $II$ and touches the center of circle $II$, the diameter of circle $I$ is the radius of circle $II$.

 Applying the area formula $A = \pi r^2$, substitute $4$ for $A$ to solve for the radius of circle $1$.
 \[4 = \pi r^2\]
 \[\frac{4}{\pi} = r^2\]
 \[r = \frac{2}{\sqrt{\pi}}\]

 That means the diameter of circle $I$ (or the radius of circle $II$) is $\frac{4}{\sqrt{\pi}}$.  Apply the area formula again to find the area of circle $II$.
 \[A = \pi (\frac{2}{\sqrt{\pi}})^2 = \pi \cdot \frac{4}{\pi} = 16\]

 The answer is $\boxed{\textbf{(D)}}$.

 


	\item (\textit{AHSME 1960}) The number $2.5252525\ldots$ can be written as a fraction. 
When reduced to lowest terms the sum of the numerator and denominator of this fraction is:


 $\textbf{(A) }7\qquad \textbf{(B) } 29\qquad \textbf{(C) }141\qquad \textbf{(D) }349\qquad \textbf{(E) }\text{none of these}$


 
 \textbf{Answer: }D



\textbf{Solution 1}

Let $x = 2.5252525\ldots$ , so
 \[100x = 252.52525\ldots\]
 \[x = 2.5252525\ldots\]
 \[99x = 250\]
 \[x = \frac{250}{99}\]

 The sum of the numerator and the denominator is $349$, so the answer is $\boxed{\textbf{(D)}}$ .

 



\textbf{Solution 2}

The number $2.5252525\ldots$ can also be written as
 \[2+\frac{5}{10} + \frac{2}{100} + \frac{5}{1000} + \frac{2}{10,000} \ldots\]This is really two infinite series -- one with first term $2$ and common difference $\frac{1}{100}$ and one with first term $\frac{5}{10}$ and common difference $\frac{1}{100}$.  Apply the infinite series formula $S = \frac{a_1}{1-r}$ on both series.
 \[\frac{2}{1-\frac{1}{100}} + \frac{\frac{5}{10}}{1-\frac{1}{100}}\]
 \[\frac{2}{\frac{99}{100}} + \frac{\frac{5}{10}}{\frac{99}{100}}\]
 \[2 \cdot \frac{100}{99} + \frac{5}{10} \cdot \frac{100}{99}\]
 \[\frac{200}{99} + \frac{50}{99}\]
 \[\frac{250}{99}\]

 The sum of the numerator and the denominator is $349$, so the answer is $\boxed{\textbf{(D)}}$ .

 


	\item (\textit{AHSME 1960}) The fraction $\frac{a^2+b^2-c^2+2ab}{a^2+c^2-b^2+2ac}$ is (with suitable restrictions of the values of a, b, and c):


 $\text{(A) irreducible}\qquad$


 $\text{(B) reducible to negative 1}\qquad$


 $\text{(C) reducible to a polynomial of three terms}\qquad$


 $\text{(D) reducible to} \frac{a-b+c}{a+b-c}\qquad$


 $\text{(E) reducible to} \frac{a+b-c}{a-b+c}$


 
 \textbf{Answer: }E



\textbf{Solution}

Use the commutative property to get
 \[\frac{a^2+2ab+b^2-c^2}{a^2+2ac+c^2-b^2}\]Factor perfect square trinomials to get
 \[\frac{(a+b)^2-c^2}{(a+c)^2-b^2}\]Factor difference of squares to get
 \[\frac{(a+b+c)(a+b-c)}{(a+b+c)(a-b+c)}\]Cancel out like terms (with suitable restrictions of a, b, c) to get
 \[\frac{a+b-c}{a-b+c}\]The answer is $\boxed{\textbf{(E)}}$.

 


	\item (\textit{AHSME 1960}) Given the following six statements:

 \[\text{(1) All women are good drivers}\]

 \[\text{(2) Some women are good drivers}\]

 \[\text{(3) No men are good drivers}\]

 \[\text{(4) All men are bad drivers}\]

 \[\text{(5) At least one man is a bad driver}\]

 \[\text{(6) All men are good drivers.}\]


 

The statement that negates statement $(6)$ is:


 

$\textbf{(A) }(1)\qquad \textbf{(B) }(2)\qquad \textbf{(C) }(3)\qquad \textbf{(D) }(4)\qquad \textbf{(E) }(5)$


 
 \textbf{Answer: }E



\textbf{Solution}

$\fbox{(E)5}$

 If you look at (5) and (6), if either one of them are false, it implies that the other one is true and vice versa≥also haiyaaaa hwy u no know solve it treeveal cheevian. @dhar

 


	\item (\textit{AHSME 1960}) For a given value of $k$ the product of the roots of $x^2-3kx+2k^2-1=0$
is $7$. The roots may be characterized as:


 $\textbf{(A) }\text{integral and positive} \qquad\textbf{(B) }\text{integral and negative} \qquad \\ \textbf{(C) }\text{rational, but not integral} \qquad\textbf{(D) }\text{irrational} \qquad\textbf{(E) } \text{imaginary}$


 
 \textbf{Answer: }D



\textbf{Solution}

If the product of the roots are $7$, then by \href{https://artofproblemsolving.com/wiki/index.php?title=Vieta\%27s\_formulas}{Vieta's formulas}, 
 \[2k^2-1=7\]Solve for $k$ in the resulting equation to get
 \[2k^2=8\]
 \[k^2=4\]
 \[k=\pm 2\]That means the two quadratics are $x^2-6x+7=0$ and $x^2+6x+7=0$.  Since $b^2$, $a$, and $c$ are the same, the \href{https://artofproblemsolving.com/wiki/index.php?title=Discriminant}{discriminant} of both is$36-(4 \cdot 1 \cdot 7) = 8$.  Because $8$ is not a perfect square, the roots for both are irrational, so the answer is $\boxed{\textbf{(D)}}$.

 


	\item (\textit{AHSME 1960}) The locus of the centers of all circles of given radius $a$, in the same plane, passing through a fixed point, is:


 $\textbf{(A) }\text{a point}\qquad \textbf{(B) }\text{ a straight line}\qquad \textbf{(C) }\text{two straight lines}\qquad \textbf{(D) }\text{a circle}\qquad  \textbf{(E) }\text{two circles}$


 
 \textbf{Answer: }D



\textbf{Solution}


                    \begin{figure}[H]
                        centering
                        \includegraphics[width=0.5\textwidth]{Solution/AHSME_I_1960_12_Solution_Figure_1.png}
                    \end{figure}
                    

 If a circle passes through a point, then the point is $a$ units away from the center.  That means that all of the centers are $a$ units from the point.  By definition, the resulting figure is a circle, so the answer is $\boxed{\textbf{(D)}}$.

 


	\item (\textit{AHSME 1960}) The polygon(s) formed by $y=3x+2, y=-3x+2$, and $y=-2$, is (are):


 $\textbf{(A) }\text{an equilateral triangle}\qquad\textbf{(B) }\text{an isosceles triangle} \qquad\textbf{(C) }\text{a right triangle} \qquad \\ \textbf{(D) }\text{a triangle and a trapezoid}\qquad\textbf{(E) }\text{a quadrilateral}$


 
 \textbf{Answer: }B



\textbf{Solution}


                    \begin{figure}[H]
                        centering
                        \includegraphics[width=0.5\textwidth]{Solution/AHSME_I_1960_13_Solution_Figure_1.png}
                    \end{figure}
                    

 The points of intersection of two of the lines are $(0,2)$ and $(\pm \frac{4}{3} , -2)$, so use the Distance Formula to find the sidelengths.

 Two of the side lengths are $\sqrt{(\frac{4}{3})^2+4^2} = \frac{4 \sqrt{10}}{3}$ while one of the side lengths is $4$.  That makes the triangle isosceles, so the answer is $\boxed{\textbf{(B)}}$.

 


	\item (\textit{AHSME 1960}) If $a$ and $b$ are real numbers, the equation $3x-5+a=bx+1$ has a unique solution $x$ [The symbol $a \neq 0$ means that $a$ is different from zero]:


 $\text{(A) for all a and b} \qquad \text{(B) if a }\neq\text{2b}\qquad \text{(C) if a }\neq 6\qquad \\ \text{(D) if b }\neq 0\qquad \text{(E) if b }\neq 3$


 
 \textbf{Answer: }E



\textbf{Solution}

If the coefficients of the x-terms are equal on both sides, then when the x-terms are subtracted from both sides, the equation results in a number equals 1.

 This means the equation has either infinite or no solutions, so the x-terms can not be equal on both sides.  Thus, $b \neq 3$, so the answer is $\boxed{\textbf{(E)}}$.

 


	\item (\textit{AHSME 1960}) Triangle $I$ is equilateral with side $A$, perimeter $P$, area $K$, and circumradius $R$ (radius of the circumscribed circle). 
Triangle $II$ is equilateral with side $a$, perimeter $p$, area $k$, and circumradius $r$. If $A$ is different from $a$, then: 


 $\textbf{(A)}\ P:p = R:r \text{ } \text{only sometimes} \qquad \textbf{(B)}\ P:p = R:r \text{ } \text{always}\qquad \\ \textbf{(C)}\ P:p = K:k \text{ } \text{only sometimes} \qquad \textbf{(D)}\ R:r = K:k \text{ } \text{always}\qquad \textbf{(E)}\ R:r = K:k \text{ } \text{only sometimes}$


 
 \textbf{Answer: }B



\textbf{Solution}


                    \begin{figure}[H]
                        centering
                        \includegraphics[width=0.5\textwidth]{Solution/AHSME_I_1960_15_Solution_Figure_1.png}
                    \end{figure}
                    

 First, find $P$, $K$, and $R$ in terms of $A$.  Since all sides of an \href{https://artofproblemsolving.com/wiki/index.php?title=Equilateral\_triangle}{equilateral triangle} are the same, $P=3A$.  From the area formula, $K=\frac{A^2\sqrt{3}}{4}$.  By using 30-60-90 triangles, $R=\frac{A\sqrt{3}}{3}$.

 Using the same steps, $p=3a$, $k=\frac{a^2\sqrt{3}}{4}$, and $r=\frac{a\sqrt{3}}{3}$.

 Note that $P/p = 3A/3a = A/a$ and $R/r = \frac{A\sqrt{3}}{3} \div \frac{a\sqrt{3}}{3} = A/a$.  That means $P/p = R/r$, so the answer is $\boxed{\textbf{(B)}}$.

 


	\item (\textit{AHSME 1960}) In the numeration system with base $5$, counting is as follows: $1, 2, 3, 4, 10, 11, 12, 13, 14, 20,\ldots$.
The number whose description in the decimal system is $69$, when described in the base $5$ system, is a number with:


 $\textbf{(A)}\ \text{two consecutive digits} \qquad\textbf{(B)}\ \text{two non-consecutive digits} \qquad \\ \textbf{(C)}\ \text{three consecutive digits} \qquad\textbf{(D)}\ \text{three non-consecutive digits} \qquad \\ \textbf{(E)}\ \text{four digits}$


 
 \textbf{Answer: }C



\textbf{Solution}

Since $25<69<125$, divide $69$ by $25$.  The quotient is $2$ and the remainder is $19$, so rewrite the number as
 \[69 = 2 \cdot 25 + 19\]Similarly, dividing $19$ by $5$ results in quotient of $3$ and remainder of $4$, so rewrite the number as
 \[69 = 2 \cdot 25 + 3 \cdot 5 + 4 \cdot 1\]Thus, the number in base $5$ can be written as $234_5$, so the answer is $\boxed{\textbf{(C)}}$

 


	\item (\textit{AHSME 1960}) The formula $N=8 \times 10^{8} \times x^{-3/2}$ gives, for a certain group, the number of individuals whose income exceeds $x$ dollars. 
The lowest income, in dollars, of the wealthiest $800$ individuals is at least:


 $\textbf{(A)}\ 10^4\qquad \textbf{(B)}\ 10^6\qquad \textbf{(C)}\ 10^8\qquad \textbf{(D)}\ 10^{12} \qquad \textbf{(E)}\ 10^{16}$


 
 \textbf{Answer: }A



\textbf{Solution}

Plug $800$ for $N$ because $800$ is the number of people who has at least $x$ dollars.
 \[800 = 8 \times 10^{8} \times x^{-3/2}\]
 \[10^{-6} = x^{-3/2}\]In order to undo raising to the $-\frac{3}{2}$ power, raise both sides to the $-\frac{2}{3}$ power.
 \[(10^{-6})^{-2/3} = (x^{-3/2})^{-2/3}\]
 \[x = 10^4  \text{ dollars}\]

 The answer is $\boxed{\textbf{(A)}}$.

 


	\item (\textit{AHSME 1960}) The pair of equations $3^{x+y}=81$ and $81^{x-y}=3$ has:


 $\textbf{(A)}\ \text{no common solution} \qquad \\ \textbf{(B)}\ \text{the solution} \text{ } x=2, y=2\qquad \\ \textbf{(C)}\ \text{the solution} \text{ } x=2\frac{1}{2}, y=1\frac{1}{2} \qquad \\ \textbf{(D)}\text{ a common solution in positive and negative integers} \qquad \\ \textbf{(E)}\ \text{none of these}$


 
 \textbf{Answer: }E



\textbf{Solution}

Rewrite the equations so both sides have a common base.
 \[3^{x+y}=3^4\]
 \[3^{4(x-y)}=3^1\]Taking the logarithm base 3, we get a linear \href{https://artofproblemsolving.com/wiki/index.php?title=System\_of\_equations}{system of equations}.
 \[x+y=4\]
 \[4x-4y=1\]Solve the system to get $x=\frac{17}{8}$ and $y=\frac{15}{8}$.  The answer is $\boxed{\textbf{(E)}}$.

 


	\item (\textit{AHSME 1960}) Consider equation $I: x+y+z=46$ where $x, y$, and $z$ are positive integers, and equation $II: x+y+z+w=46$, 
where $x, y, z$, and $w$ are positive integers. Then


 $\textbf{(A)}\ \text{I can be solved in consecutive integers} \qquad \\ \textbf{(B)}\ \text{I can be solved in consecutive even integers} \qquad \\ \textbf{(C)}\ \text{II can be solved in consecutive integers} \qquad \\ \textbf{(D)}\ \text{II can be solved in consecutive even integers} \qquad \\ \textbf{(E)}\ \text{II can be solved in consecutive odd integers}$


 
 \textbf{Answer: }C



\textbf{Solution}

Consider each option, one at a time.

 For option A, let $x=y-1$ and $z=y+1$.  That means $3y=46$, so $y=\frac{46}{3}$.  That is not an integer, so option A is eliminated.

 For option B, let $x=y-2$ and $z=y+2$.  That also means $3y=46$, so $y=\frac{46}{3}$.  That is also not an integer, so option B is eliminated.

 For option C, let $y=x+1$, $z=x+2$, and $w=x+3$.  That means $4x+6=46$, so $x=10$.  Option C works.

 For options D and E, let $y=x+2$, $z=x+4$, and $w=x+6$.  That means $4x+12=46$, so $x=\frac{17}{2}$.  Since the result is not an integer, options D and E are eliminated.

 Thus, the answer is $\boxed{\textbf{(C)}}$.

 


	\item (\textit{AHSME 1960}) The coefficient of $x^7$ in the expansion of $(\frac{x^2}{2}-\frac{2}{x})^8$ is:


 $\textbf{(A)}\ 56\qquad \textbf{(B)}\ -56\qquad \textbf{(C)}\ 14\qquad \textbf{(D)}\ -14\qquad \textbf{(E)}\ 0$


 
 \textbf{Answer: }D



\textbf{Solution}

By the \href{https://artofproblemsolving.com/wiki/index.php?title=Binomial\_Theorem}{Binomial Theorem}, each term of the expansion is $\binom{8}{n}\left(\frac{x^2}{2}\right)^{8-n}\left(\frac{-2}{x}\right)^n$.

 We want the exponent of $x$ to be $7$, so 
 \[2(8-n)-n=7\]
 \[16-3n=7\]
 \[n=3\]

 If $n=3$, then the corresponding term is
 \[\binom{8}{3}\left(\frac{x^2}{2}\right)^{5}\left(\frac{-2}{x}\right)^3\]
 \[56 \cdot \frac{x^{10}}{32} \cdot \frac{-8}{x^3}\]
 \[-14x^7\]

 The answer is $\boxed{\textbf{(D)}}$.

 


	\item (\textit{AHSME 1960}) The diagonal of square $I$ is $a+b$. The perimeter of square $II$ with twice the area of $I$ is:


 $\textbf{(A)}\ (a+b)^2\qquad \textbf{(B)}\ \sqrt{2}(a+b)^2\qquad \textbf{(C)}\ 2(a+b)\qquad \textbf{(D)}\ \sqrt{8}(a+b) \qquad \textbf{(E)}\ 4(a+b)$


 
 \textbf{Answer: }E



\textbf{Solution}

Since the diagonal of square $I$ is $a+b$ units long, the side length of square $I$ is $\frac{a+b}{\sqrt{2}}$, so the area of square $I$ is $\frac{(a+b)^2}{2}$.

 The area of square $II$ is twice as much as the area of square $I$, so the area of square $II$ is $(a+b)^2$.  That means the side length of square $II$ is $a+b$, so the perimeter of square $II$ is $4(a+b)$, or answer choice $\boxed{\textbf{(E)}}$.

 


	\item (\textit{AHSME 1960}) The equality $(x+m)^2-(x+n)^2=(m-n)^2$, where $m$ and $n$ are unequal non-zero constants, is satisfied by $x=am+bn$, where:


 $\textbf{(A)}\ a = 0, b \text{ } \text{has a unique non-zero value}\qquad \\ \textbf{(B)}\ a = 0, b \text{ } \text{has two non-zero values}\qquad \\ \textbf{(C)}\ b = 0, a \text{ } \text{has a unique non-zero value}\qquad \\ \textbf{(D)}\ b = 0, a \text{ } \text{has two non-zero values}\qquad \\ \textbf{(E)}\ a \text{ } \text{and} \text{ } b \text{ } \text{each have a unique non-zero value}$


 
 \textbf{Answer: }A



\textbf{Solution}

Expand binomials, combine like terms, and subtract terms from both sides.
 \[x^2 + 2xm + m^2 - x^2 - 2xn - n^2 = m^2 - 2mn + n^2\]
 \[2xm + m^2 - 2xn - n^2 = m^2 - 2mn + n^2\]
 \[2xm - 2xn - n^2 = -2mn + n^2\]Get all the x-terms on one side and factor to solve for x.
 \[2xm - 2xn = -2mn + 2n^2\]
 \[2x(m-n) = -2n(m-n)\]Since $m \not= n$, both sides can be divided by $m-n$.
 \[x = -n\]That means $a = 0$ and $b = -1$, so the answer is $\boxed{\textbf{(A)}}$.

 


	\item (\textit{AHSME 1960}) The radius $R$ of a cylindrical box is $8$ inches, the height $H$ is $3$ inches. 
The volume $V = \pi R^2H$ is to be increased by the same fixed positive amount when $R$ 
is increased by $x$ inches as when $H$ is increased by $x$ inches. This condition is satisfied by:


 $\textbf{(A)}\ \text{no real value of x} \qquad \\ \textbf{(B)}\ \text{one integral value of x} \qquad \\ \textbf{(C)}\ \text{one rational, but not integral, value of x} \qquad \\ \textbf{(D)}\ \text{one irrational value of x}\qquad \\ \textbf{(E)}\ \text{two real values of x}$


 
 \textbf{Answer: }C



\textbf{Solution}

Since increasing the height by $x$ inches should result in the same volume as increasing the radius by $x$ inches, write an equation with the two cylinders (one with height increased, one with radius increased).
 \[\pi (8+x)^2 \cdot 3 = \pi \cdot 8^2 \cdot (3+x)\]
 \[3(x^2+16x+64) = 64(x+3)\]
 \[3x^2+48x+192=64x+192\]
 \[3x^2-16x=0\]
 \[x(3x-16)=0\]By the Zero-Product Property, $x = 0$ or $\frac{16}{3}$.  However, since $x$ must be increased, discard $x = 0$ as a possible value.  Thus, the length should be increased by $\frac{16}{3}$ inches, so the answer is $\boxed{\textbf{(C)}}$.

 


	\item (\textit{AHSME 1960}) If $\log_{2x}216 = x$, where $x$ is real, then $x$ is:


 $\textbf{(A)}\ \text{A non-square, non-cube integer}\qquad$
$\textbf{(B)}\ \text{A non-square, non-cube, non-integral rational number}\qquad$
$\textbf{(C)}\ \text{An irrational number}\qquad$
$\textbf{(D)}\ \text{A perfect square}\qquad$
$\textbf{(E)}\ \text{A perfect cube}$


 
 \textbf{Answer: }A



\textbf{Solution 1}

Rewrite the equation to $(2x)^x = 216$.  Since $2x$ is the base of the logarithm, we only need to check x-values that are positive.

 With trial and error, $3$ is a solution because $(2 \cdot 3)^3 = 6^3 = 216$.  Since $(2x)^x$ gets larger as x gets larger, $3$ is the only solution, so the answer is $\boxed{\textbf{(A)}}$.

 


 


 


 


 


 


 


 


 


 


 


 


 



\textbf{Solution 2}

We use the definition of logarithms: if $\log_b a = x$, then $b^x = a$. Applying this to the given equation $\log_{2x} 216 = x$, we obtain$(2x)^x = 216$.

 Rewrite $216$ as $6^3$. By checking small positive integer values, we find that $x = 3$ satisfies the equation since$(2 \cdot 3)^3 = 6^3 = 216$.

 Thus, $x = 3$, which is a non-square, non-cube integer. Hence, the correct answer is $\boxed{\text{A}}$.$\text{(Solution by Trashboycan)}$

 
.

 


	\item (\textit{AHSME 1960}) Let $m$ and $n$ be any two odd numbers, with $n$ less than $m$. 
The largest integer which divides all possible numbers of the form $m^2-n^2$ is:


 $\textbf{(A)}\ 2\qquad \textbf{(B)}\ 4\qquad \textbf{(C)}\ 6\qquad \textbf{(D)}\ 8\qquad \textbf{(E)}\ 16$


 
 \textbf{Answer: }D



\textbf{Solution}

First, factor the \href{https://artofproblemsolving.com/wiki/index.php?title=Difference\_of\_squares}{difference of squares}.
 \[(m+n)(m-n)\]Since $m$ and $n$ are odd numbers, let $m=2a+1$ and $n=2b+1$, where $a$ and $b$ can be any integer.
 \[(2a+2b+2)(2a-2b)\]Factor the resulting expression.
 \[4(a+b+1)(a-b)\]If $a$ and $b$ are both even, then $a-b$ is even.  If $a$ and $b$ are both odd, then $a-b$ is even as well.  If $a$ is odd and $b$ is even (or vise versa), then $a+b+1$ is even.  Therefore, in all cases, $8$ can be divided into all numbers with the form $m^2-n^2$.

 This can be confirmed by setting $m=3$ and $n=1$, making $m^2-n^2=9-1=8$.  Since $8$ is not a multiple of $3$ and is less than $16$, we can confirm that the answer is $\boxed{\textbf{(D)}}$.

 \textbf{solution 2(logic)} we can factorize it into 
 \[(m+n)(m-n)\] the we knows both of them are even and if we have 2 consecutive even number atlesast one of the me is divisible by 4 so our answer is 4*2 =8

 


	\item (\textit{AHSME 1960}) Find the set of $x$-values satisfying the inequality $\left|\frac{5-x}{3}\right| < 2$. [The symbol $|a|$ means $+a$ if $a$ is positive, 
$-a$ if $a$ is negative, $0$ if $a$ is zero. The notation $1 < a < 2$ means that $a$ can have any value between $1$ and $2$, excluding $1$ and $2$.]


 $\textbf{(A)}\ 1 < x < 11\qquad \textbf{(B)}\ -1 < x < 11\qquad \textbf{(C)}\ x < 11\qquad \textbf{(D)}\ x > 11\qquad \textbf{(E)}\ |x| < 6$


 
 \textbf{Answer: }B



\textbf{Solution 1}

Break up the \href{https://artofproblemsolving.com/wiki/index.php?title=Absolute\_value}{absolute value} into two cases.

 For the first case, let $x < 5$, so $\frac{5-x}{3}$ is positive.  That means (for $x<5$)
 \[\frac{5-x}{3} < 2\]
 \[5-x<6\]
 \[-x<1\]
 \[x>-1\]For the second case, let $x \ge 5$, so $\frac{5-x}{3}$ is negative.  That means (for $x \ge 5$)
 \[\frac{x-5}{3} < 2\]
 \[x-5<6\]
 \[x<11\]

 Combine both cases to get $-1 < x < 11$, which is answer choice $\boxed{\textbf{(B)}}$.

 



\textbf{Solution 2}


                    \begin{figure}[H]
                        centering
                        \includegraphics[width=0.5\textwidth]{Solution/AHSME_I_1960_26_Solution_Figure_1.png}
                    \end{figure}
                    

 Another way to solve this is to graph $y=|\frac{5-x}{3}|$ and $y=2$.  The solution is the areas on the graph where the y-values of $y=|\frac{5-x}{3}|$ are lower than $2$.  From the graph, $-1<x<11$, so the answer is $\boxed{\textbf{(B)}}$.

 


	\item (\textit{AHSME 1960}) Let $S$ be the sum of the interior angles of a polygon $P$ for which each interior angle is $7\frac{1}{2}$ times the 
exterior angle at the same vertex. Then


 $\textbf{(A)}\ S=2660^{\circ} \text{ } \text{and} \text{ } P \text{ } \text{may be regular}\qquad \\ \textbf{(B)}\ S=2660^{\circ} \text{ } \text{and} \text{ } P \text{ } \text{is not regular}\qquad \\ \textbf{(C)}\ S=2700^{\circ} \text{ } \text{and} \text{ } P \text{ } \text{is regular}\qquad \\ \textbf{(D)}\ S=2700^{\circ} \text{ } \text{and} \text{ } P \text{ } \text{is not regular}\qquad \\ \textbf{(E)}\ S=2700^{\circ} \text{ } \text{and} \text{ } P \text{ } \text{may or may not be regular}$


 
 \textbf{Answer: }E



\textbf{Solution}

Let $a_n$ be the interior angle of the nth vertex, and let $b_n$ be the exterior angle of the nth vertex.  From the conditions in the problem,
 \[a_n = 7.5b_n\]That means
 \[a_1 + a_2 \cdots a_n = 7.5(b_1 + b_2 \cdots b_n)\]Since the sum of the exterior angles of a polygon is $360^{\circ}$, the equation can be simplified as
 \[a_1 + a_2 \cdots a_n = 7.5 \cdot 360\]
 \[a_1 + a_2 \cdots a_n = 2700\]The sum of the interior angles is $2700$ degrees.  However, there is no other constraint on what angle the interior angles can be, so we can not for sure claim that the polygon is regular or not.  Thus, the answer is $\boxed{\textbf{(E)}}$.

 


 


	\item (\textit{AHSME 1960}) The equation $x-\frac{7}{x-3}=3-\frac{7}{x-3}$ has:


 $\textbf{(A)}\ \text{infinitely many integral roots}\qquad\textbf{(B)}\ \text{no root}\qquad\textbf{(C)}\ \text{one integral root}\qquad$
$\textbf{(D)}\ \text{two equal integral roots}\qquad\textbf{(E)}\ \text{two equal non-integral roots}$


 
 \textbf{Answer: }B



\textbf{Solution}

Both terms have a $-\frac{7}{x-3}$ term, so add $\frac{7}{x-3}$ to both sides.  This results in $x = 3$.

 However, note that if $3$ is plugged back into the original equation, it results in
 \[3-\frac{7}{0}=3-\frac{7}{0}\]

 Since dividing by zero is undefined, $3$ is an extraneous solution.  That means there are no solutions, so the answer is $\boxed{\textbf{(B)}}$.

 


	\item (\textit{AHSME 1960}) Five times $A$'s money added to $B$'s money is more than $\$51.00$. Three times $A$'s money minus $B$'s money is $\$21.00$.
If $a$ represents $A$'s money in dollars and $b$ represents $B$'s money in dollars, then:


 $\textbf{(A)}\ a>9, b>6 \qquad \textbf{(B)}\ a>9, b<6 \qquad \textbf{(C)}\ a>9, b=6\qquad \textbf{(D)}\ a>9, \text{but we can put no bounds on} \text{ } b\qquad \textbf{(E)}\ 2a=3b$


 
 \textbf{Answer: }A



\textbf{Solution}

Use math symbols to write an equation and an inequality based on the conditions in the problem.
 \[5a+b>51\]
 \[3a-b=21\]From the second equation, $b = 3a-21$ and $a = 7 + b/3$.

 Substituting $b$ in the inequality results in
 \[5a + 3a - 21 > 51\]
 \[8a > 72\]
 \[a > 9\]

 Substituting $a$ in the inequality results in
 \[35 + \frac{5b}{3} + b > 51\]
 \[35 + \frac{8b}{3} > 51\]
 \[\frac{8b}{3} > 16\]
 \[b > 6\]

 The answer is $\boxed{\textbf{(A)}}$.

 


	\item (\textit{AHSME 1960}) Given the line $3x+5y=15$ and a point on this line equidistant from the coordinate axes. Such a point exists in:


 $\textbf{(A)}\ \text{none of the quadrants}\qquad\textbf{(B)}\ \text{quadrant I only}\qquad\textbf{(C)}\ \text{quadrants I, II only}\qquad$
$\textbf{(D)}\ \text{quadrants I, II, III only}\qquad\textbf{(E)}\ \text{each of the quadrants}$


 
 \textbf{Answer: }C



\textbf{Solution}

If a point is equidistant from the coordinate axes, then the absolute values of the x-coordinate and y-coordinate are equal.  Since the point is on the line $3x+5y=15$, find the intersection point of $y=x$ and $3x+5y=15$ and the intersection point of $y=-x$ and $3x+5y=15$.

 Substituting $x$ for $y$ in $3x+5y=15$ results in $8x=15$, so $x=\frac{15}{8}$.  That means $y=\frac{15}{8}$, so one of the points is in the first quadrant.

 Substituting $-x$ and $y$ in $3x+5y=15$ results in $-2x=15$, so $x=\frac{-15}{2}$.  That means $y=\frac{15}{2}$, so the other point is in the second quadrant.

 Thus, the points are in quadrants $I$ and $II$, so the answer is $\boxed{\textbf{(C)}}$.

 


	\item (\textit{AHSME 1960}) For $x^2+2x+5$ to be a factor of $x^4+px^2+q$, the values of $p$ and $q$ must be, respectively:


 $\textbf{(A)}\ -2, 5\qquad \textbf{(B)}\ 5, 25\qquad \textbf{(C)}\ 10, 20\qquad \textbf{(D)}\ 6, 25\qquad \textbf{(E)}\ 14, 25$


 
 \textbf{Answer: }D



\textbf{Solution}

Let the other quadratic be $x^2 + bx + c$, where $(x^2 + 2x + 5)(x^2 + bx + c) = x^4 + px^2 + q$.  Multiply the two quadratics to get
 \[x^4 + (b+2)x^3 + (c + 2b + 5)x^2 + (2c + 5b)x + 5c\]Since $x^4 + px^2 + q$ have no $x^3$ term and no $x$ term, the coefficients of these terms must be zero.
 \[b+2=0\]
 \[b=-2\]
 \[2c+5b=0\]
 \[2c-10=0\]
 \[c=5\]Substitute $c$ and $b$ back to get
 \[x^4 + 6x^2 + 25\]Thus, $p=6$ and $q=25$, so the answer is $\boxed{\textbf{(D)}}$.

 


	\item (\textit{AHSME 1960}) In this figure the center of the circle is $O$. $AB \perp BC$, $ADOE$ is a straight line, $AP = AD$, and $AB$ has a length twice the radius. Then:


 
                    \begin{figure}[H]
                        centering
                        \includegraphics[width=0.5\textwidth]{Problem/AHSME_I_1960_32_Problem_Figure_1.png}
                    \end{figure}
                    


 $\textbf{(A)} AP^2 = PB \times AB\qquad \\ \textbf{(B)}\ AP \times DO = PB \times AD\qquad \\ \textbf{(C)}\ AB^2 = AD \times DE\qquad \\ \textbf{(D)}\ AB \times AD = OB \times AO\qquad \\ \textbf{(E)}\ \text{none of these}$


 
 \textbf{Answer: }A



\textbf{Solution 1}

We claim that $\fbox{A}$ is the right answer.

 Let the radius of circle $O$ be $r$, and let the length of $AD=x$. Since $AB \perp BC$, $AB$ is a tangent to circle $O$. Thus, by the tangent-secant theorem, we have $AB^2=AD\times AE$, or, $(2r)^2=x(x+2r)$. Through some algebraic manipulation, we find 
 \begin{align}AD^2=x^2=2r(2r-x)\end{align}. 

 Since $AD^2=AP^2=x^2$, $PB=AB-AP=2r-x$, and $AB=2r$, we see that $(1)$ is identical to $AP^2=PB\times AB$, hence our answer is $\fbox{A}$

 



\textbf{Solution 2 (guess and check)}

Let $r$ be the radius of the circle, so $AB = 2r$.  By the \href{https://artofproblemsolving.com/wiki/index.php?title=Pythagorean\_Theorem}{Pythagorean Theorem}, $AO = \sqrt{(2r)^2 + r^2} = r \sqrt{5}$.  That means, $AD = AP = r \sqrt{5} - r$, so $PB = 2r - r\sqrt{5} + r = 3r - r\sqrt{5}$.

 Substitute values for each answer choice to determine which one is correct for all $r$.

 For option A, substitution results in
 \[(r \sqrt{5} - r)^2 = (3r - r\sqrt{5})2r\]
 \[5r^2 - 2r^2 \sqrt{5} + r^2 = 6r^2 - 2r^2 \sqrt{5}\]
 \[6r^2 - 2r^2 \sqrt{5} = 6r^2 - 2r^2 \sqrt{5}\]

 For option B, substitution results in
 \[(r \sqrt{5} - r)r = (3r - r\sqrt{5})(r \sqrt{5} - r)\]
 \[r^2 \sqrt{5} - r^2 = 4r^2 \sqrt{5} - 8r^2\]

 For option C, substitution results in
 \[(2r)^2 = (r\sqrt{5} - r)2r\]
 \[4r^2 = 2r^2 \sqrt{5} - 2r^2\]

 For option D, substitution results in
 \[(r \sqrt{5} - r) \cdot 2r = r \cdot r\sqrt{5}\]
 \[2r^2 \sqrt{5} - 2r^2 = r^2 \sqrt{5}\]

 From each option, only option A has both sides equaling each other, so the answer is $\boxed{\textbf{(A)}}$.

 


	\item (\textit{AHSME 1960}) You are given a sequence of $58$ terms; each term has the form $P+n$ where $P$ stands for the product $2 \times 3 \times 5 \times\ldots \times 61$
of all prime numbers less than or equal to $61$, and $n$ takes, successively, the values $2, 3, 4,\ldots, 59$. 
Let $N$ be the number of primes appearing in this sequence. Then $N$ is:


 $\textbf{(A)}\ 0\qquad \textbf{(B)}\ 16\qquad \textbf{(C)}\ 17\qquad \textbf{(D)}\ 57\qquad \textbf{(E)}\ 58$


 
 \textbf{Answer: }A



\textbf{Solution}

First, note that $n$ does not have a \href{https://artofproblemsolving.com/wiki/index.php?title=Prime\_number}{prime number} larger than $61$ as one of its factors.  Also, note that $n$ does not equal $1$.

 Therefore, since the prime factorization of $n$ only has primes from $2$ to $59$, $n$ and $P$ share at least one common factor other than $1$.  Therefore $P+n$ is not prime for any $n$, so the answer is $\Rightarrow{\boxed{\textbf{(A)}}}$.

 


	\item (\textit{AHSME 1960}) Two swimmers, at opposite ends of a $90$-foot pool, start to swim the length of the pool, 
one at the rate of $3$ feet per second, the other at $2$ feet per second. 
They swim back and forth for $12$ minutes. Allowing no loss of times at the turns, find the number of times they pass each other. 


 $\textbf{(A)}\ 24\qquad \textbf{(B)}\ 21\qquad \textbf{(C)}\ 20\qquad \textbf{(D)}\ 19\qquad \textbf{(E)}\ 18$


 
 \textbf{Answer: }C



\textbf{Solution}

First, note that it will take $30$ seconds for the first swimmer to reach the other side and $45$ seconds for the second swimmer to reach the other side.  Also, note that after $180$ seconds (or $3$ minutes), both swimmers will complete an even number of laps, essentially returning to their starting point.

 
                    \begin{figure}[H]
                        centering
                        \includegraphics[width=0.5\textwidth]{Solution/AHSME_I_1960_34_Solution_Figure_1.png}
                    \end{figure}
                    

 At this point, find the number of meeting points in the first $3$ minutes, then multiply by four to get the answer.  From the graph (where the x-axis is the time in seconds and the y-axis is distance from one side of the pool), there are five meeting points, so the two swimmers will pass each other $20$ times, which is answer choice $\boxed{\textbf{(C)}}$.

 


	\item (\textit{AHSME 1960}) From point $P$ outside a circle, with a circumference of $10$ units, a tangent is drawn. 
Also from $P$ a secant is drawn dividing the circle into unequal arcs with lengths $m$ and $n$. 
It is found that $t_1$, the length of the tangent, is the mean proportional between $m$ and $n$. 
If $m$ and $t$ are integers, then $t$ may have the following number of values:


 $\textbf{(A)}\ \text{zero}\qquad \textbf{(B)}\ \text{one}\qquad \textbf{(C)}\ \text{two}\qquad \textbf{(D)}\ \text{three}\qquad \textbf{(E)}\ \text{infinitely many}$


 
 \textbf{Answer: }C



\textbf{Solution}

By definition of \href{https://artofproblemsolving.com/wiki/index.php?title=Geometric\_mean}{mean proportional}, $t = \sqrt{mn}$.  Since $m+n=10$, $t = \sqrt{m(10-m)}$.

 With trial and error, note that when $m=9$, $t=3$ and when $m=2$, $t=4$.  These values work since another \href{https://artofproblemsolving.com/wiki/index.php?title=Tangent\_line}{tangent line} can be drawn from $P$, and the angle between the \href{https://artofproblemsolving.com/wiki/index.php?title=Tangent\_line}{tangent} and \href{https://artofproblemsolving.com/wiki/index.php?title=Secant\_line}{secant} can decrease to match the values of $m$ and $n$.  Thus, the answer is $\boxed{\textbf{(C)}}$.

 


	\item (\textit{AHSME 1960}) Let $s_1, s_2, s_3$ be the respective sums of $n, 2n, 3n$ terms of the same arithmetic progression with $a$ as the first term and $d$ 
as the common difference. Let $R=s_3-s_2-s_1$. Then $R$ is dependent on:


 $\textbf{(A)}\ a\text{ }\text{and}\text{ }d\qquad \textbf{(B)}\ d\text{ }\text{and}\text{ }n\qquad \textbf{(C)}\ a\text{ }\text{and}\text{ }n\qquad \textbf{(D)}\ a, d,\text{ }\text{and}\text{ }n\qquad  \textbf{(E)}\ \text{neither} \text{ } a \text{ } \text{nor} \text{ } d \text{ } \text{nor} \text{ } n$


 
 \textbf{Answer: }B



\textbf{Solution}

The nth term of the sequence with first term $a$ is $a + d(n-1)$.  That means the sum of the first $n$ terms is with first term $a$ is $\frac{n(2a + dn -d)}{2}$.

 The 2nth term of the sequence with first term $a$ is $a + d(2n-1)$.  That means the sum of the first $2n$ terms is with first term $a$ is $\frac{2n(2a + 2dn -d)}{2}$.

 The 3nth term of the sequence with first term $a$ is $a + d(3n-1)$.  That means the sum of the first $3n$ terms is with first term $a$ is $\frac{3n(2a + 3dn -d)}{2}$.

 Substituting $s_1$, $s_2$, and $s_3$ with its respective values results in

 
 \[\frac{3n(2a + 3dn -d)}{2}-\frac{2n(2a + 2dn -d)}{2}-\frac{n(2a + dn -d)}{2}\]
 \[\frac{6an + 9dn^2 -3nd}{2}-\frac{4an + 4dn^2 - 2nd}{2}-\frac{2an + dn^2 - dn}{2}\]

 
 \[\frac{4dn^2}{2}\]

 
 \[2dn^2\]

 The answer is $\boxed{\textbf{(B)}}$.

 


 


	\item (\textit{AHSME 1960}) The base of a triangle is of length $b$, and the altitude is of length $h$. 
A rectangle of height $x$ is inscribed in the triangle with the base of the rectangle in the base of the triangle. The area of the rectangle is:


 $\textbf{(A)}\ \frac{bx}{h}(h-x)\qquad \textbf{(B)}\ \frac{hx}{b}(b-x)\qquad \textbf{(C)}\ \frac{bx}{h}(h-2x)\qquad \textbf{(D)}\ x(b-x)\qquad \textbf{(E)}\ x(h-x)$


 
 \textbf{Answer: }A



\textbf{Solution}

Let $AB=b$, $DE=h$, and $WX = YZ = x$.
                    \begin{figure}[H]
                        centering
                        \includegraphics[width=0.5\textwidth]{Solution/AHSME_I_1960_37_Solution_Figure_1.png}
                    \end{figure}
                    Since $CD$ is perpendicular to $AB$, $ND = WX$.  That means $CN = h-x$.The sides of the rectangle are parallel, so $XY \parallel WZ$.  That means by AA Similarity, $\triangle CXY \sim \triangle CAB$.  Letting $n$ be the length of the base of the rectangle, that means
 \[\frac{h-x}{n} = \frac{h}{b}\]
 \[n = \frac{b(h-x)}{h}\]Thus, the area of the rectangle is $\frac{bx}{h}(h-x)$, which is answer choice $\boxed{\textbf{(A)}}$.

 


 


	\item (\textit{AHSME 1960}) In this diagram $AB$ and $AC$ are the equal sides of an isosceles $\triangle ABC$, in which is inscribed equilateral $\triangle DEF$. 
Designate $\angle BFD$ by $a$, $\angle ADE$ by $b$, and $\angle FEC$ by $c$. Then:


 
                    \begin{figure}[H]
                        centering
                        \includegraphics[width=0.5\textwidth]{Problem/AHSME_I_1960_38_Problem_Figure_1.png}
                    \end{figure}
                    


 $\textbf{(A)}\ b=\frac{a+c}{2}\qquad \textbf{(B)}\ b=\frac{a-c}{2}\qquad \textbf{(C)}\ a=\frac{b-c}{2} \qquad \textbf{(D)}\ a=\frac{b+c}{2}\qquad \textbf{(E)}\ \text{none of these}$


 
 \textbf{Answer: }D



\textbf{Solution}

Since $\triangle DEF$ is an \href{https://artofproblemsolving.com/wiki/index.php?title=Equilateral\_triangle}{equilateral triangle}, all of the angles are $60^{\circ}$.The angles in a line add up to $180^{\circ}$, so
 \[\angle FDB = 120 - b\]
 \[\angle EFC = 120 - a\]The angles in a triangle add up to $180^{\circ}$, so
 \[\angle ABC = 60 + b - a\]
 \[\angle ACB = 60 - c + a\]Since $\triangle ABC$ is \href{https://artofproblemsolving.com/wiki/index.php?title=Isosceles\_triangle}{isosceles} and $AB = AC$, by Base-Angle Theorem,
 \[60 + b - a = 60 - c + a\]
 \[b + c = 2a\]
 \[a = \frac{b+c}{2}\]The answer is $\boxed{\textbf{(D)}}$.

 


 


	\item (\textit{AHSME 1960}) To satisfy the equation $\frac{a+b}{a}=\frac{b}{a+b}$, $a$ and $b$ must be:


 $\textbf{(A)}\ \text{both rational}\qquad\textbf{(B)}\ \text{both real but not rational}\qquad\textbf{(C)}\ \text{both not real}\qquad$
$\textbf{(D)}\ \text{one real, one not real}\qquad\textbf{(E)}\ \text{one real, one not real or both not real}$


 



 
 \textbf{Answer: }E



\textbf{Solution}

First, note that $a \neq 0$ and $a \neq -b$.  Cross multiply both sides to get
 \[a^2 + 2ab + b^2 = ab\]Subtract both sides by $ab$ to get
 \[a^2 + ab + b^2 = 0\]From the \href{https://artofproblemsolving.com/wiki/index.php?title=Quadratic\_formula}{quadratic formula},
 \[a = \frac{-b \pm \sqrt{b^2 - 4b^2}}{2}\]
 \[a = \frac{-b \pm \sqrt{-3b^2}}{2}\]If $b$ is \href{https://artofproblemsolving.com/wiki/index.php?title=Real}{real}, then $\sqrt{-3b^2}$ is imaginary because $-3b^2$ is negative.If $b$ is not real, where $b = m+ni$ and $n \neq 0$, then $\sqrt{-3b^2}$ evaluates to $\sqrt{-3m^2 - 6mni + 3n^2}$.  As long as $m \neq 0$, the expression can also be imaginary because a real number squared will be a real number.From these two points, the answer is $\boxed{\textbf{(E)}}$.

 


	\item (\textit{AHSME 1960}) Given right $\triangle ABC$ with legs $BC=3, AC=4$. Find the length of the shorter angle trisector from $C$ to the hypotenuse:


 $\textbf{(A)}\ \frac{32\sqrt{3}-24}{13}\qquad\textbf{(B)}\ \frac{12\sqrt{3}-9}{13}\qquad\textbf{(C)}\ 6\sqrt{3}-8\qquad\textbf{(D)}\ \frac{5\sqrt{10}}{6}\qquad\textbf{(E)}\ \frac{25}{12}\qquad$


 
 \textbf{Answer: }A



\textbf{Solution}

Angle $C$ is split into three $30^{\circ}$ angles. The shorter angle trisector will be the one closer $BC$. Let it intersect $AB$ at point $P$. Let the perpendicular from point $P$ intersect $BC$ at point $R$ and have length $x$. Thus $\triangle PRC$ is a $30^{\circ}-60^{\circ}-90^{\circ}$ triangle and $RC$ has length $x\sqrt{3}$. Because $\triangle PBR$ is similar to $\triangle ABC$, $RB$ has length $\frac{3}{4}x$. 
 \[RC+RB=BC=x\sqrt{3}+\frac{3}{4}x=3\]The problem asks for the length of $PC$, or $2x$. Solving for $x$ and multiplying by two gives $\boxed{\textbf(A)}$.

 


\end{enumerate}